\scanpage{101}
\mspage{\unsure{95|45}}
\diagram{\begin{tikzcd}[ampersand replacement=\&,row sep=large,column sep=large]
E'_{j,n}\arrow[r,no head,"\unsure{v_i|u_i}"]\&E^{\deleted{\prime}}_{i,n}\\
E'_{k,n}\arrow[u]\arrow[r,no head,"\unsure{v_j|u_j}"]\&E_{j,n}\arrow[u]
\end{tikzcd}}
où les flèches verticales sont les homomorphismes de transition.

N.B. On pose $E_{i,n}=E_i/I_i^{n+1}$, $E'_{i,n}=\operatorname{Symm}_{C_i}(I_i/I_i^2)/$ puissance $(N+1)$ème de l'\unsure{idéan|idéal} d'augmentation. La démonstration de lemme B est immédiate, en procédant comme dans le cas particulier 2.2. déjà traité.

\uline{Lemme} C Supposons $C$ et $B$ form simple sur $A$, (donc en vertu de lemme A les conditions de Lemme B sont vérifiées). Alors:

\inserted{a)} Pour tout \deleted{$i$, \unsure{existe un} $j\unsure{>|,}i$ et un $v_i$ comme dans lemme B a), avec en plus un} système $i,j,v_i$ comme dans Lemme B\inserted{a)}, \deleted{il existe un}\inserted{il existe} $k\geq j$ \inserted{et}\deleted{\unsure{qui soit} rendant valable la conclusion de lemme B, b), et tel qu'il existe} un homomorphisme de $A$-algèbres $w_j:E_{k,n}\mathbin{-}E'_{j,n}$ ayant les propriétés suivantes: $1^\circ$) \inserted{$w_{\unsure{n|k}}$} induit un épimorphisme sur les gradués associés $2^\circ$) $v_iw_j$ est l'homomorphisme de transition $E_k\to E_{\unsure{i}}$

\deleted{b) Soient $i,j,k,v_i,v_j,w_j$ comme dans a). Alors il existe un $\unsure{\ell}\quad k$, un homomorphisme $v_k:E'_{\unsure{\ell|k},n}\quad E_{k,n}$ et un homomorphisme $w_k:E_{\unsure{\ell|k},n}\ \deleted{\unsure{E}}\ E'_{k,n}$ ayant les propriétés suivantes: $1^\circ$) $v_k$ a les propriétés mentionnées déjà pour $v_i,v_j$ $2^\circ$) $w_k$ a les deux propriétés mentionnées déjà pour $w_j$ $3^\circ$}
\margin{100}
